\section{Representación matemática y respuesta esperada sin controlador}
\label{sec:p5}
\enunciado{De una representación matemática del sistema según la teoría. Explique cómo debería
quedar la respuesta final de la planta (sin controlador) y algunos posibles valores que se puedan
extraer teóricamente}

\subsection{Representación matemática del sistema}
El sistema de frenado se modela como un disco en rotación que, después de alcanzar una velocidad
angular inicial $\omega_0$, deja de recibir torque de accionamiento del motor DC y desacelera por
efecto de la fricción. La planta tiene como entrada la acción del freno $u(t)$, aplicada mediante
una goma movida por un servomotor MG90S, y como salida la velocidad angular $\omega(t)$ del disco.
Según la segunda ley de Newton para sistemas rotacionales, la suma de torques es igual al momento
de inercia equivalente multiplicado por la aceleración angular~\cite{openstax_rotacion}:
\begin{equation}
  J_\text{eq}\,\frac{d\omega(t)}{dt} = \tau_m(t) - b\,\omega(t) - \tau_f(t)
  \label{eq:newton_rot}
\end{equation}
En esta expresión, $J_\text{eq}$ representa el momento de inercia equivalente
(\si{kg\cdot m^2}), $\tau_m$ el torque del motor (\si{N\cdot m}), $b$ el coeficiente de fricción
viscosa (\si{N\cdot m\cdot s/rad}) y $\tau_f$ el torque resistivo del freno (\si{N\cdot m}).
Durante la etapa de frenado se aproxima $\tau_m \approx 0$, aunque los elementos que permanecen
acoplados al eje siguen contribuyendo a la inercia y a las pérdidas mecánicas. Para obtener una
función de transferencia aproximada, se consideran variaciones pequeñas $\Delta\omega$ y
$\Delta u$ y se supone localmente $\Delta\tau_f \approx K_f\,\Delta u$. Sustituyendo en la
ecuación rotacional se obtiene la ecuación incremental, y posteriormente se aplica Laplace con
condiciones iniciales incrementales nulas:
\begin{gather*}
  J_\text{eq}\,\frac{d\,\Delta\omega(t)}{dt} + b\,\Delta\omega(t) = -K_f\,\Delta u(t) \\[4pt]
  (J_\text{eq}s + b)\,\Delta\Omega(s) = -K_f\,\Delta U(s)
\end{gather*}
\begin{equation}
  G(s) = \frac{\Delta\Omega(s)}{\Delta U(s)} \approx \frac{-K_f}{J_\text{eq}s + b}
  \label{eq:planta_G}
\end{equation}
El signo negativo indica que incrementar el frenado disminuye la velocidad. Esta función de
transferencia es una aproximación lineal alrededor de un intervalo de operación, no un modelo
válido para toda la detención, pues la fricción, la presión de la goma y el comportamiento del
servo pueden ser no lineales~\cite{ctms_motor}. Esta simplificación se adopta porque permite
estudiar la relación entre la acción de frenado y la velocidad mediante la transformada de
Laplace, y proporciona un punto de partida para el diseño del controlador PID. La validez de la
aproximación no se da por supuesta: posteriormente se comprobará mediante la comparación del
modelo con datos experimentales en MATLAB, delimitando las condiciones de operación en las que
resulte representativo.

\subsection{Respuesta esperada de la planta sin controlador}
Sin PID, el servomotor aplica una posición de frenado fija y no corrige los cambios de velocidad
medidos. Si durante esa prueba se aproxima el torque resistivo total como constante y positivo,
$\tau_\text{res}$, la ecuación de movimiento se simplifica a:
\begin{equation}
  J_\text{eq}\,\frac{d\omega(t)}{dt} = -\tau_\text{res}
  \label{eq:torque_constante}
\end{equation}
Despejando la derivada de la velocidad e integrando con la condición inicial
$\omega(0)=\omega_0$, se obtiene la respuesta temporal. Se impone además $\omega(t)\geq 0$, ya
que la acción del freno no invierte el sentido del movimiento:
\begin{gather*}
  \frac{d\omega}{dt} = -\frac{\tau_\text{res}}{J_\text{eq}}
  \quad\Longrightarrow\quad
  \omega(t) = \omega_0 - \frac{\tau_\text{res}}{J_\text{eq}}\,t
\end{gather*}
\begin{equation}
  \omega(t) = \max\left\{0,\ \omega_0 - \frac{\tau_\text{res}}{J_\text{eq}}\,t\right\}
  \label{eq:respuesta_lineal}
\end{equation}
Por lo tanto, se espera una curva de velocidad aproximadamente lineal y decreciente hasta llegar a
\SI{0}{rad/s}. Una mayor acción de frenado produciría una pendiente más negativa y un menor tiempo
de detención. En el prototipo real pueden observarse desviaciones respecto a esa recta debido al
rozamiento viscoso, la fricción seca y el retardo de aplicación del freno. Incluso sin ejercer
presión con la zapata, el disco eventualmente se detendría por las pérdidas mecánicas; por eso se
deben comparar las curvas con y sin frenado aplicado.

\subsection{Valores que pueden extraerse teóricamente}
Las magnitudes teóricas se expresan en función de los parámetros del sistema, sin asignarles
valores numéricos. En general, el momento de inercia de un cuerpo que gira alrededor de un eje se
calcula a partir de su distribución de masa. Para un disco macizo y homogéneo, de masa $m$ y radio
$R$, esta relación se simplifica a~\cite{openstax_rotacion}:
\begin{equation}
  J_\text{disco} = \int r^2\,dm
  \quad\Longrightarrow\quad
  J_\text{disco} = \frac{1}{2}\,mR^2
  \label{eq:inercia_disco}
\end{equation}
Si varios elementos giran alrededor del mismo eje, su inercia se incorpora al momento equivalente
mediante $J_\text{eq} = J_\text{disco} + \sum J_i$. Bajo la hipótesis de torque resistivo total
constante, es posible determinar simbólicamente la aceleración angular y el tiempo de detención:
\begin{equation}
  \alpha = -\frac{\tau_\text{res}}{J_\text{eq}}\,;\qquad
  t_\text{det} = \frac{J_\text{eq}\,\omega_0}{\tau_\text{res}}
  \label{eq:alfa_tdet}
\end{equation}
La velocidad inicial se expresa en rad/s mediante $\omega_0 = 2\pi\cdot\text{RPM}_0/60$. Las
ecuaciones de movimiento son generales para sistemas rotacionales; la expresión
$J_\text{disco} = \tfrac{1}{2}mR^2$ corresponde específicamente a un disco uniforme. En esta etapa
no se asignan valores reales a los parámetros: posteriormente se identificarán y ajustarán en
MATLAB utilizando mediciones experimentales de entrada y velocidad, y se validará la capacidad
predictiva del modelo.
